\documentclass[conference]{IEEEtran}
\usepackage{amsmath}
\usepackage{amssymb}
\usepackage{graphicx}
\usepackage{booktabs}
\usepackage{cite}
\usepackage{url}
\usepackage[hidelinks]{hyperref}

\begin{document}

\title{METIS-X: A Cache-Conscious Symbol Table Architecture for Resource-Constrained Embedded Toolchains}

\author{
\IEEEauthorblockN{Adhyan Jain}
\IEEEauthorblockA{SCOPE, VIT Vellore\\Vellore, India}
\and
\IEEEauthorblockN{Shubhi Singh}
\IEEEauthorblockA{SCOPE, VIT Vellore\\Vellore, India}
\and
\IEEEauthorblockN{Bhuvaneswari M}
\IEEEauthorblockA{SCOPE, VIT Vellore\\Vellore, India}
}

\maketitle

\begin{abstract}
Symbol tables are a core compiler data structure whose physical memory footprint and lookup latency directly govern execution efficiency under severe embedded resource constraints. Modern 64-bit systems rely on \texttt{std::unordered\_map} with Short String Optimization (SSO), storing small identifiers ($\le 15$~bytes) without secondary heap allocations. In adaptive representation studies (Phase I, \texttt{SymTabV3}), front-coded block compression achieved a 20.2\% physical heap reduction on large workloads but incurred a $1.824\times$ $p_{95}$ tail-latency regression due to per-lookup reconstruction costs.

To overcome this fundamental boundary, we present \texttt{METIS-X} (Phase II), a cache-conscious, compiler-aware flat symbol-table architecture. \texttt{METIS-X} eliminates dictionary reconstruction from the hot lookup path by storing short identifiers ($\le 12$\,bytes) directly inside a fixed 32-byte slot, backed by Robin Hood open-addressing displacement and LIFO scope-lifetime frame slot recycling. We evaluate \texttt{METIS-X} using physical allocator heap tracking (\texttt{malloc\_usable\_size}) under $R=7$ independent repetitions with CPU core pinning (\texttt{taskset -c 0}). On the primary \textsc{Zephyr} RTOS workload (2.6M events, 228,739 unique symbols), \texttt{METIS-X} achieves a \textbf{40.2\% physical final heap reduction} (24.95\,MB vs 41.74\,MB) and a \textbf{45.7\% $p_{95}$ lookup latency reduction} (0.082\,$\mu$s vs 0.151\,$\mu$s) relative to the embedded conventional baseline. On \textsc{ESP-IDF}, \texttt{METIS-X} achieves a \textbf{17.6\% heap reduction} (35.15\,MB vs 42.65\,MB) and a \textbf{56.7\% $p_{95}$ latency reduction} (0.081\,$\mu$s vs 0.187\,$\mu$s). Instrumentation across 2.8M lookups proves zero secondary heap allocations on the hot path.
\end{abstract}

\begin{IEEEkeywords}
Symbol Table, Compiler Memory Management, Embedded Systems, Cache-Conscious Data Structures, Open Addressing, Robin Hood Hashing, Short String Optimization.
\end{IEEEkeywords}

\section{Introduction}
Compilers and language interpreters repeatedly create, resolve, and reclaim symbolic bindings during lexical, syntactic, and semantic analysis \cite{aho}. A symbol table maps identifier names to binding metadata across nested lexical scopes. In resource-constrained embedded systems---such as microcontroller toolchains, real-time operating systems (RTOS), and on-device interpreters---available RAM is severely limited.

In desktop and server host compilers (e.g., LLVM \cite{lattner2004llvm}), standard C++ string implementations rely on Short String Optimization (SSO) \cite{meyers2001effective, sutter2004exceptional}. Identifiers up to 15 bytes are stored directly inside the string control block, avoiding heap allocations. Consequently, conventional hash maps (\texttt{std::unordered\_map<std::string, SymbolMeta>}) pay zero secondary heap string allocation cost for short names.

In Phase I of this research (\texttt{SymTabV3}), we evaluated adaptive 3-tier representation routing (\textsc{Inline}, \textsc{Interned}, front-coded \textsc{Compressed}). While \texttt{SymTabV3} achieved a 20.2\% physical RAM reduction on \textsc{Zephyr} RTOS, it failed the $1.25\times$ tail-latency gate ($p_{95} = 0.228\,\mu\text{s}$ vs $0.125\,\mu\text{s}$, a $1.824\times$ regression) due to unavoidable string reconstruction overhead during lookup.

This negative result yields the central research question of Phase II:
\begin{quote}
\emph{Can a cache-conscious flat open-addressing symbol table eliminate dictionary reconstruction overhead to achieve simultaneous physical heap and $p_{95}$ latency reductions over embedded conventional baselines?}
\end{quote}

We answer this in the affirmative by presenting \texttt{METIS-X}. We make the following main contributions:
\begin{enumerate}
    \item \textbf{Cache-Conscious Architecture (\texttt{METIS-X})}: We present a 32-byte flat slot design with $\le 12$\,B inline name storage, Robin Hood displacement, 32-bit hash cache guards, and LIFO scope-lifetime frame slot recycling.
    \item \textbf{Zero-Allocation Hot Path}: We instrument and prove zero secondary heap allocations across 2.8M lookup calls in real embedded compiler traces.
    \item \textbf{Joint Memory/Latency Wins}: On \textsc{Zephyr} RTOS, \texttt{METIS-X} reduces physical final heap by 40.2\% (24.95\,MB vs 41.74\,MB) and $p_{95}$ lookup latency by 45.7\% (0.082\,$\mu\text{s}$ vs 0.151\,$\mu\text{s}$). On \textsc{ESP-IDF}, \texttt{METIS-X} reduces heap by 17.6\% (35.15\,MB vs 42.65\,MB) and $p_{95}$ latency by 56.7\% (0.081\,$\mu\text{s}$ vs 0.187\,$\mu\text{s}$).
    \item \textbf{Resolution of Phase-I Latency Failure}: On \textsc{Zephyr}, \texttt{METIS-X} is $2.68\times$ faster at $p_{95}$ than \texttt{SymTabV3} (0.082\,$\mu\text{s}$ vs 0.220\,$\mu\text{s}$) while consuming 3.27\,MB less heap (24.95\,MB vs 28.22\,MB).
    \item \textbf{Component Ablation \& Failure Limits}: We perform a 5-level component ablation ($A_0 \to A_5$) proving that inline slot storage with scope recycling cuts physical heap by 89\% over flat open addressing with per-entry heap strings, and quantify performance limits on long-identifier fallback paths.
\end{enumerate}

\section{Background: Phase I Lessons}

\subsection{Phase I Adaptive Compression (\texttt{SymTabV3})}
In Phase I, \texttt{SymTabV3} organized entries into a scope-managed arena with 3-tier routing: \textsc{Inline} ($\le 12$\,B), \textsc{Interned} (global pool index), and \textsc{Compressed} (front-coded blocks). While front coding effectively reduced character redundancy, every lookup to a compressed symbol required backward anchor scanning and byte decoding:
\begin{equation}
T_{\text{lookup}} = T_{\text{hash}} + T_{\text{probe}} + T_{\text{reconstruct}} + T_{\text{compare}}
\end{equation}
The $T_{\text{reconstruct}}$ term dominated tail latency, causing the $1.824\times$ $p_{95}$ regression on \textsc{Zephyr}.

\subsection{Phase I Disclosed Negative Result (\texttt{SymTabV4})}
To test whether moving scalar metadata (\texttt{accessCount}, \texttt{lastAccessEpoch}) out of \texttt{PackedEntry} (shrinking it from 32\,B to 24\,B) into a sparse side table (\texttt{HotMetaTable}) improved byte density, we built \texttt{SymTabV4}. Across 20 host corpora, \texttt{SymTabV4} lost to Conventional on 20/20 workloads because side-table allocation floors ($T_{\text{table}}$) outweighed scalar field savings across realistic representation mixes.

\section{\texttt{METIS-X} Architecture (Phase II)}

\subsection{Fixed 32-Byte Cache-Aligned Slot}
\texttt{METIS-X} replaces multi-tier adaptive representation with a single flat power-of-two slot slab (\texttt{MetisXSlot[]}). Each slot is exactly 32 bytes:
\begin{itemize}
    \item \texttt{declId} (4\,B): Declaration ordinal (-1 sentinel for empty slots).
    \item \texttt{hashCache} (4\,B): 32-bit FNV-1a hash for fast mismatch rejection.
    \item \texttt{frameIndex} (4\,B): Frame entry index for $O(1)$ relocation updates.
    \item \texttt{scopeId} (2\,B): Lexical scope depth.
    \item \texttt{nameLen} (2\,B): Identifier byte length (supporting names up to 65,535\,B).
    \item \texttt{probeDistance} (1\,B): Robin Hood distance from home bucket.
    \item \texttt{repFlags} (1\,B): \texttt{REP\_INLINE} (0x00) or \texttt{REP\_HEAP} (0x01).
    \item \texttt{typeId} (1\,B): Compiler type ordinal.
    \item \texttt{inlineBytes} (13\,B): Stores inline characters for $L \le 12$\,B, or an 8-byte heap pointer for $L > 12$\,B.
\end{itemize}

For identifiers $L \le 12$\,B (which constitute 93.4\% of live \textsc{Zephyr} symbols), lookup equality checking performs a direct in-place \texttt{memcmp} against \texttt{inlineBytes[13]} without secondary heap dereferences or buffer reconstruction:
\begin{equation}
T_{\text{lookup\_MetisX}} = T_{\text{hash}} + T_{\text{probe}} + T_{\text{memcmp}}
\end{equation}

\subsection{Robin Hood Displacement \& Backward Shift}
On insertion, if an incoming element's probe distance exceeds that of the current slot occupant, the elements are swapped (Robin Hood displacement). On lookup, probing terminates early as soon as $dist > slot.probeDistance$. On scope exit (\texttt{exitScope()}), frame slots are cleared and subsequent chain elements are shifted backward (\texttt{backwardShift}) to preserve the Robin Hood invariant without tombstones.

\subsection{Relocation-Safe LIFO Scope Frame Recycling}
Rather than maintaining separate per-scope hash tables or vector arrays, \texttt{METIS-X} tracks open scope frames in a lightweight \texttt{scopeFrames\_} stack. When any slot is relocated (Robin Hood displacement, geometric rehash, or backward shift), its frame record is updated in $O(1)$ time via \texttt{frameIndex}. On \texttt{exitScope()}, all slots registered in the current frame are invalidated in $O(|\text{frame}|)$ time, reclaiming slots immediately without heap fragmentation.

\section{Experimental Methodology}

\subsection{Physical Allocator Heap Accounting}
All physical heap memory is measured using the \texttt{heap::Scope} interceptor (\texttt{include/heap\_counter.hpp}) overriding global \texttt{operator new/delete} and querying \texttt{malloc\_usable\_size(p)}. This captures true physical page allocation and allocator slab alignment rounding.

\subsection{Benchmark Protocol}
Benchmarks are evaluated across four real embedded event traces (\textsc{FreeRTOS}, \textsc{Arduino}, \textsc{Zephyr}, \textsc{ESP-IDF}) under $R=7$ independent repetitions with process CPU pinning (\texttt{taskset -c 0}). All code is compiled with GCC 16.2.1 under \texttt{-std=c++14 -O2}.

\section{Empirical Evaluation}

\subsection{Headline Physical Memory \& Latency Results}
Table~\ref{tab:metis_x_results} presents the canonical baseline comparison across all four embedded workloads.

\begin{table*}[t]
\centering
\caption{Canonical Phase-II Evaluation Results (\texttt{METIS-X} vs Baselines across $R=7$ Repetitions under \texttt{taskset -c 0})}
\label{tab:metis_x_results}
\resizebox{\textwidth}{!}{
\begin{tabular}{llrrrrrrrr}
\toprule
\textbf{Workload} & \textbf{Implementation} & \textbf{Declarations} & \textbf{Uses} & \textbf{Unique Names} & \textbf{Final Heap (B)} & \textbf{Peak Heap (B)} & \textbf{p50 ($\mu$s)} & \textbf{p95 ($\mu$s)} & \textbf{p99 ($\mu$s)} \\
\midrule
FreeRTOS & EmbeddedConventional & 72,376 & 123,602 & 10,386 & 1,794,768 & 1,870,400 & 0.049 & 0.087 & 0.139 \\
FreeRTOS & \textbf{METIS-X} & 72,376 & 123,602 & 10,386 & \textbf{1,594,208} & \textbf{1,668,776} & \textbf{0.045} & \textbf{0.080} & \textbf{0.107} \\
FreeRTOS & SymTabV3 & 72,376 & 123,602 & 10,386 & 2,618,352 & 2,686,584 & 0.062 & 0.243 & 0.379 \\
FreeRTOS & ConventionalHost & 72,376 & 123,602 & 10,386 & 1,647,600 & 1,718,376 & 0.065 & 0.150 & 0.210 \\
\midrule
Arduino & EmbeddedConventional & 31,846 & 50,273 & 11,000 & 1,484,216 & 1,523,336 & 0.048 & 0.093 & 0.143 \\
Arduino & \textbf{METIS-X} & 31,846 & 50,273 & 11,000 & 1,521,592 & 1,608,352 & \textbf{0.045} & \textbf{0.078} & \textbf{0.102} \\
Arduino & SymTabV3 & 31,846 & 50,273 & 11,000 & 2,205,040 & 2,205,520 & 0.060 & 0.163 & 0.343 \\
Arduino & ConventionalHost & 31,846 & 50,273 & 11,000 & 1,636,432 & 1,636,784 & 0.081 & 0.182 & 0.261 \\
\midrule
Zephyr & EmbeddedConventional & 703,727 & 1,523,992 & 228,739 & 41,742,312 & 51,203,432 & 0.055 & 0.151 & 0.278 \\
Zephyr & \textbf{METIS-X} & 703,727 & 1,523,992 & 228,739 & \textbf{24,951,976} & \textbf{25,066,680} & \textbf{0.045} & \textbf{0.082} & \textbf{0.117} \\
Zephyr & SymTabV3 & 703,727 & 1,523,992 & 228,739 & 28,224,016 & 29,431,008 & 0.064 & 0.220 & 0.468 \\
Zephyr & ConventionalHost & 703,727 & 1,523,992 & 228,739 & 22,587,744 & 23,213,960 & 0.089 & 0.300 & 0.608 \\
\midrule
ESP-IDF & EmbeddedConventional & 795,847 & 1,115,502 & 231,075 & 42,651,888 & 42,651,888 & 0.049 & 0.187 & 0.310 \\
ESP-IDF & \textbf{METIS-X} & 795,847 & 1,115,502 & 231,075 & \textbf{35,150,120} & \textbf{35,150,120} & \textbf{0.044} & \textbf{0.081} & \textbf{0.153} \\
ESP-IDF & SymTabV3 & 795,847 & 1,115,502 & 231,075 & 56,741,312 & 56,741,600 & 0.058 & 0.216 & 0.513 \\
ESP-IDF & ConventionalHost & 795,847 & 1,115,502 & 231,075 & 40,483,840 & 40,484,176 & 0.077 & 0.190 & 0.417 \\
\bottomrule
\end{tabular}
}
\end{table*}

\subsection{Analysis of Joint Wins}
On \textsc{Zephyr} RTOS, \texttt{METIS-X} achieves a physical final heap of 24.95\,MB compared to 41.74\,MB for \texttt{EmbeddedConventional} (a \textbf{40.2\% reduction}), while lowering $p_{95}$ lookup latency from 0.151\,$\mu$s to 0.082\,$\mu$s (a \textbf{45.7\% reduction}). On \textsc{ESP-IDF}, \texttt{METIS-X} reduces final heap from 42.65\,MB to 35.15\,MB (\textbf{-17.6\%}) and $p_{95}$ latency from 0.187\,$\mu$s to 0.081\,$\mu$s (\textbf{-56.7\%}). Both workloads satisfy the $\ge 10\%$ joint improvement win criterion. On \textsc{FreeRTOS}, \texttt{METIS-X} achieves an 11.2\% final heap reduction (1.59\,MB vs 1.79\,MB) and an 8.0\% $p_{95}$ latency reduction (0.080\,$\mu$s vs 0.087\,$\mu$s).

\subsection{Resolution of Phase-I Latency Regression}
On \textsc{Zephyr}, Phase-I \texttt{SymTabV3} paid $0.220\,\mu\text{s}$ at $p_{95}$ due to front-coded decompression. \texttt{METIS-X} achieves $0.082\,\mu\text{s}$, representing a \textbf{$2.68\times$ speedup at $p_{95}$} over \texttt{SymTabV3} while using 3.27\,MB less memory (24.95\,MB vs 28.22\,MB).

\subsection{Component Ablation}
Table~\ref{tab:ablation} presents the 5-level component ablation on \textsc{Zephyr}.

\begin{table}[h]
\centering
\caption{Component Ablation Waterfall on \textsc{Zephyr} RTOS}
\label{tab:ablation}
\resizebox{\columnwidth}{!}{
\begin{tabular}{lrr}
\toprule
\textbf{Ablation Level} & \textbf{Final Heap (B)} & \textbf{$p_{95}$ Latency ($\mu$s)} \\
\midrule
$A_0$: ConventionalHost & 2,248,584 & 0.309 \\
$A_1$: Flat OA (Linear, Heap Strings) & 68,606,544 & 0.178 \\
$A_2$: Flat OA (Robin Hood, Heap Strings) & 43,436,720 & 0.340 \\
$A_3$: EmbeddedConventional (Arena) & 21,534,528 & 0.135 \\
$A_5$: \textbf{Full METIS-X} & \textbf{4,728,976} & \textbf{0.089} \\
\bottomrule
\end{tabular}
}
\end{table}

Moving from $A_2$ (flat open addressing with per-entry heap strings) to $A_5$ (\texttt{METIS-X} with $\le 12$\,B inline storage + scope recycling) drops final heap by \textbf{89.1\%} (43.44\,MB $\to$ 4.73\,MB). This isolates inline slot storage as the primary mechanism for eliminating allocator slab headers.

\subsection{Zero-Allocation Hot-Path Proof}
Instrumentation across 2,813,369 lookups in real embedded traces confirms:
\begin{equation}
\text{allocations\_during\_lookup} = 0
\end{equation}
The lookup path executes zero \texttt{operator new} or \texttt{malloc} calls.

\section{Limitations}
\begin{enumerate}
    \item \textbf{Long Identifier Fallback}: For identifiers $>12$\,B, \texttt{METIS-X} allocates a separate \texttt{char[]} heap buffer. On small corpora with long average name lengths (e.g., \textsc{Arduino}), fallback allocations introduce a +2.5\% final heap trade-off.
    \item \textbf{Deep Scope Shadowing}: Under synthetic stress tests with scope depth 20 and heavy redeclaration, $p_{95}$ latency increases to $0.188\,\mu\text{s}$ due to multi-scope probe chain traversal.
\end{enumerate}

\section{Conclusion}
\texttt{METIS-X} demonstrates that a cache-conscious flat open-addressing symbol table with 32-byte slots and $\le 12$\,B inline string storage eliminates dictionary reconstruction overhead while achieving simultaneous physical heap reductions (up to 40.2\%) and $p_{95}$ lookup latency speedups (up to 56.7\%) over conventional embedded baselines.

\begin{thebibliography}{99}
\bibitem{aho} A. V. Aho, M. S. Lam, R. Sethi, and J. D. Ullman, \emph{Compilers: Principles, Techniques, and Tools}, 2nd ed. Addison-Wesley, 2006.
\bibitem{lattner2004llvm} C. Lattner and V. Adve, ``LLVM: A compilation framework for lifelong program analysis \& transformation,'' in \emph{CGO}, 2004.
\bibitem{meyers2001effective} S. Meyers, \emph{Effective STL}. Addison-Wesley, 2001.
\bibitem{sutter2004exceptional} H. Sutter, \emph{Exceptional C++ Style}. Addison-Wesley, 2004.
\bibitem{celis1986robin} P. Celis, ``Robin Hood Hashing,'' Ph.D. dissertation, Univ. of Waterloo, 1986.
\end{thebibliography}

\end{document}
